\chapter{Avaliação}

O processo de avaliação é fundamental para assegurar a qualidade da formação acadêmica, orientar o aperfeiçoamento das práticas pedagógicas e garantir que os objetivos do curso sejam atingidos de forma coerente e efetiva. A avaliação ocorre em dois níveis complementares: o da aprendizagem discente e o do curso como um todo.

\section{Avaliação da Aprendizagem}

A avaliação da aprendizagem no curso de \ppccurso é formativa, contínua, diagnóstica e somativa, conforme previsto no Regimento Geral da UFU e nas diretrizes pedagógicas da FEELT. Seu principal objetivo é acompanhar o desenvolvimento progressivo das competências cognitivas, técnicas e atitudinais previstas no perfil do egresso.

As principais características desse processo são:

\begin{itemize}
    \item Caráter processual, com instrumentos variados aplicados ao longo do semestre letivo;
    \item Foco em competências, privilegiando a capacidade do(a) estudante em aplicar conhecimentos em situações reais ou simuladas;
    \item Diversidade de métodos, incluindo provas escritas e orais, relatórios técnicos, projetos, atividades práticas, autoavaliações, apresentações, seminários, experimentações e resolução de problemas;
    \item Critérios transparentes, previamente definidos nos planos de ensino, permitindo que o(a) estudante conheça os objetivos e os parâmetros de avaliação desde o início do componente curricular;
    \item Oportunidade de recuperação, assegurada conforme o Regulamento da Graduação da UFU, para estudantes que não alcançarem o rendimento mínimo satisfatório;
    \item Registro sistemático, realizado por meio de sistema de gestão acadêmica, com acesso garantido ao(à) discente.
\end{itemize}

A avaliação não se limita à aferição de resultados, mas constitui um instrumento de acompanhamento, feedback e reorientação pedagógica, em consonância com os princípios da aprendizagem significativa, da interdisciplinaridade e do ensino centrado no estudante.

\section{Avaliação do Curso}

A avaliação do curso de Engenharia de Computação da UFU ocorre de forma permanente, participativa e orientada para a melhoria contínua, envolvendo docentes, discentes, técnico-administrativos e coordenação.

Esse processo se dá por meio de:

\begin{itemize}
    \item Autoavaliação institucional coordenada pela CPA (Comissão Própria de Avaliação), com aplicação de instrumentos periódicos junto à comunidade acadêmica;
    \item Relatórios semestrais de acompanhamento da coordenação do curso, com base em indicadores de rendimento, evasão, trancamentos, participação em programas acadêmicos e empregabilidade dos egressos;
    \item Revisão e atualização do Projeto Pedagógico do Curso, realizada periodicamente com base nos dados da autoavaliação, demandas da sociedade e atualizações das diretrizes nacionais;
    \item Participação ativa do NDE (Núcleo Docente Estruturante), responsável por analisar o desenvolvimento do currículo e propor ajustes em componentes curriculares, metodologias e estratégias avaliativas;
    \item Mecanismos de escuta ativa e representação discente, por meio do colegiado do curso e da participação dos estudantes em comissões internas;
    \item Instrumentos externos de avaliação, como o ENADE, cujos resultados são analisados e utilizados na formulação de ações de melhoria do desempenho coletivo e da formação.
\end{itemize}
    
A Avaliação do Curso é parte integrante do processo de gestão acadêmica participativa e visa promover uma cultura institucional de autocrítica, inovação e excelência na formação de engenheiros(as) de computação.